\documentclass[runningheads]{llncs}

\usepackage[T1]{fontenc}
\usepackage{graphicx}
\usepackage{booktabs}
\usepackage{multirow}
\usepackage{amsmath}
\usepackage{url}
\usepackage{xcolor}

\begin{document}

\title{GRU-Assisted Performance Enhancing Proxy for Adaptive Transmission over Dynamic High-Latency Networks}

\titlerunning{GRU-Assisted PEP for Dynamic High-Latency Networks}

\author{Guangshuo Dong\inst{1}\orcidID{0000-0000-0000-0000} \and
Pengjin Xie\inst{1}\orcidID{0000-0000-0000-0000}}

\authorrunning{G. Dong and P. Xie}

\institute{Beijing University of Posts and Telecommunications, Beijing, China\\
\email{dgs20040108@bupt.edu.cn}\quad \email{xiepengjin@bupt.edu.cn}}

\maketitle

\begin{abstract}
Stable transmission quality over dynamic high-latency networks is difficult to maintain because end-to-end control suffers from long feedback loops, delayed link-state perception, and slow recovery after link changes. Performance Enhancing Proxies (PEPs) mitigate this problem by moving buffering, local control, and recovery functions closer to bottleneck links, but conventional proxy systems may still react slowly when bandwidth and loss conditions vary. This paper proposes a GRU-assisted performance enhancing proxy for adaptive transmission over dynamic high-latency networks. Built on PEPesc, the proposed framework combines hybrid residual GRU-based link perception, an ACCEL/HOLD/BRAKE three-state decision mechanism, conservative pacing correction, and an FEC reliability floor. Experiments show that the method improves receiver throughput by about 7.6\% across five static scenarios, reduces recovery time from 1.089~s to 0.664~s in a representative dynamic scenario, decreases throughput deficit area by about 20.0\%, lowers p95 queue delay by about 68.0\%, and increases effective running time under long-duration dynamic switching by about 86.7\%.
\keywords{Performance Enhancing Proxy \and High-Latency Networks \and Dynamic Links \and GRU-Assisted Perception \and Forward Error Correction \and Adaptive Transmission}
\end{abstract}

\section{Introduction}
\label{sec:introduction}

High-latency networks are common in satellite communication, non-terrestrial networking, intercontinental transmission, and remote access systems. Compared with conventional terrestrial links, such environments often exhibit long propagation delay, high bandwidth-delay product, variable bandwidth, and non-congestion losses. End-to-end TCP control relies on acknowledgements, round-trip time (RTT) samples, and loss events to adjust its sending behavior. When the RTT is large, however, the feedback loop becomes too slow to track rapid link variations. During the delayed response interval, the sender or proxy may continue injecting traffic according to an outdated state, which can lead to queue buildup, throughput fluctuation, and slow recovery after link switching.

Performance Enhancing Proxies (PEPs) are an engineering approach to mitigate link-related performance degradation by placing local processing functions closer to the bottleneck link. By splitting part of the control loop and introducing local buffering, pacing, and recovery mechanisms, PEPs can reduce the dependence on end-to-end feedback. PEPesc extends this idea by combining a distributed TCP performance enhancing proxy with streaming forward error correction (FEC), which reduces the need for RTT-scale retransmission after losses. Nevertheless, the original PEPesc design is more suitable for static or slowly changing links. Under dynamic bandwidth and loss conditions, the proxy-side estimator may lag behind the reference link state, the control action lacks explicit state semantics, and pacing control is not sufficiently coordinated with FEC redundancy.

This paper focuses on improving the quality of networking service for dynamic high-latency transmission without redesigning a complete end-to-end congestion control protocol. Instead, we enhance the PEPesc proxy-side loop with a perception-decision-execution framework. At the perception layer, we keep the rule-based estimator as a stable baseline and introduce a hybrid residual GRU model to correct bandwidth and loss-rate estimation errors. At the decision layer, we design an ACCEL/HOLD/BRAKE state machine to distinguish link recovery, stable transmission, and congestion suppression. At the execution layer, we apply conservative pacing correction and an FEC reliability floor to avoid unstable adaptation and preserve basic streaming recovery capability.

The main contributions of this paper are summarized as follows:
\begin{itemize}
    \item We propose a GRU-assisted link perception method for PEP-side dynamic link estimation. The method learns residual corrections over a rule-based baseline rather than replacing the baseline estimator entirely.
    \item We design a three-state adaptive control mechanism with explicit ACCEL, HOLD, and BRAKE semantics, enabling interpretable control decisions under dynamic bandwidth, queue, RTT, ACK, and loss conditions.
    \item We implement a conservative execution strategy that combines legacy baseline behavior with adaptive correction and an FEC reliability floor, reducing the risk that adaptive control weakens streaming recovery.
    \item We evaluate the enhanced system using Mininet-based high-latency emulation. The results show improvements in receiver throughput, recovery time, throughput deficit area, p95 queue delay, and long-duration effective running time.
\end{itemize}

The remainder of this paper is organized as follows. Section~\ref{sec:related} reviews related work. Section~\ref{sec:framework} presents the proposed framework. Section~\ref{sec:setup} describes the implementation and experimental setup. Section~\ref{sec:evaluation} reports the evaluation results. Section~\ref{sec:discussion} discusses the findings and limitations, and Section~\ref{sec:conclusion} concludes the paper.

\section{Background and Related Work}
\label{sec:related}

\subsection{High-Latency Transmission and End-to-End Control}
High-latency transmission has long been studied in satellite and other high-bandwidth-delay-product networks. Standard TCP mechanisms can be extended with window scaling, selective acknowledgements, timestamps, and congestion-control improvements, but the sender still depends on delayed end-to-end feedback. Model-based and delay-aware algorithms such as BBR and Copa improve the trade-off between throughput and delay by estimating path properties or optimizing utility objectives. Learning-based congestion control further explores online adaptation under changing environments. These studies provide important reference points, but their control decisions are still made at the endpoint. In dynamic high-RTT links, the feedback delay can remain a fundamental obstacle when bandwidth and loss conditions switch within a time scale comparable to or shorter than the RTT.

\subsection{Performance Enhancing Proxies}
PEPs were proposed to mitigate link-related performance degradations by introducing intermediate processing such as connection splitting, local acknowledgement, caching, compression, and protocol conversion. In high-latency networks, PEPs are valuable because they move part of the control and recovery process closer to the bottleneck link. PEPesc is the direct baseline of this work. It deploys distributed proxy entities around the high-latency link and combines proxy-side bandwidth estimation with streaming coding recovery. This design reduces the dependence on RTT-scale retransmission and provides a practical basis for proxy-enhanced high-latency transmission. However, when bandwidth and loss rate vary rapidly, the original rule-based perception and control logic may not react quickly enough, leaving room for dynamic perception and control enhancement.

\subsection{Explicit Feedback and Local Control}
Network-side explicit feedback approaches, including XCP, RCP, and ABC, shorten the control loop by allowing bottleneck-side or network-side elements to participate in rate adjustment. ABC is particularly relevant because it abstracts control behavior into accelerate and brake semantics under time-varying wireless links. Such methods inspire the design of interpretable control states. Nevertheless, they typically require support from routers or network devices. In contrast, this paper assumes that only the proxy nodes can be controlled. Therefore, we construct a proxy-internal state machine that mimics the clarity of explicit feedback while relying only on locally observable proxy-side signals.

\subsection{Streaming FEC and Loss Recovery}
In high-RTT networks, loss recovery through retransmission incurs at least one additional feedback loop and can significantly delay application-layer delivery. FEC and on-the-fly erasure coding reduce this dependence by injecting redundancy before losses are fully known. PEPesc adopts streaming FEC to recover losses continuously during transmission. This mechanism is especially important when random loss is present. However, redundancy control must be coordinated with pacing: too little redundancy may interrupt decoding and feedback, whereas excessive redundancy consumes bottleneck bandwidth and can worsen queue buildup. This trade-off motivates the FEC reliability floor in our execution layer.

\subsection{Learning-Assisted Network Perception}
Network states such as bandwidth, RTT, ACK timing, queue delay, and loss evolve over time and exhibit temporal dependency. Recurrent neural networks, especially GRU and LSTM variants, are suitable for modeling such time series. In system and networking tasks, however, learning components must be deployed conservatively because inaccurate predictions can destabilize the control loop. This paper therefore uses GRU only at the perception layer. The model learns residual corrections over a rule-based estimator, while the final control decisions remain constrained by explicit states, conservative pacing correction, and FEC protection.

\section{GRU-Assisted PEP Framework}
\label{sec:framework}

\subsection{Design Overview}
The proposed framework enhances PEPesc through three coordinated layers: perception, decision, and execution. The perception layer estimates dynamic link bandwidth and loss conditions from proxy-side observations. The decision layer maps the perceived state and feedback trends into explicit control states. The execution layer applies bounded pacing correction and FEC protection to avoid unstable behavior. The framework is designed around two principles. First, the original PEPesc logic should remain the stable baseline, so the learning component only corrects estimation errors rather than directly replacing the control loop. Second, dynamic adaptation should improve recovery and queue control without degrading steady-state transmission capability.

\subsection{Hybrid Residual GRU Perception}
The perception layer receives observations collected at the proxy, including estimated bandwidth, probing bandwidth, RTT, queue delay, ACK interval, loss rate, and their temporal variations. A rule-based estimator is preserved as the baseline because it provides deterministic behavior and robustness under abnormal inputs. However, rule-based estimation can lag behind link changes near switching boundaries. To improve dynamic perception, we introduce a hybrid residual GRU model.

Let $\hat{b}_{r}$ and $\hat{l}_{r}$ denote the bandwidth and loss-rate estimates produced by the rule-based estimator. The GRU model predicts residual corrections $\Delta b$ and $\Delta l$ from a recent observation window. The final perceived states are computed as
\begin{equation}
\hat{b}=\max(0,\hat{b}_{r}+\Delta b), \quad
\hat{l}=\max(0,\hat{l}_{r}+\Delta l),
\label{eq:residual}
\end{equation}
where $\hat{b}$ and $\hat{l}$ are used by the subsequent decision logic. This residual design reduces the burden on the neural model: it does not need to learn the entire estimation process from scratch, but only learns the systematic deviation between rule-based outputs and reference link states. Non-negative constraints and smoothing are applied to suppress short-term prediction noise.

\subsection{ACCEL/HOLD/BRAKE Decision Mechanism}
Dynamic link control benefits from explicit control semantics. We divide the decision output into three states. \textbf{ACCEL} indicates that the link appears to be recovering or underutilized, so pacing may be moderately increased. \textbf{HOLD} indicates that the current transmission state is relatively stable, so the system should avoid unnecessary perturbation. \textbf{BRAKE} indicates that congestion or feedback deterioration is likely, so pacing should be suppressed to reduce queue pressure.

The state is determined by combining perceived bandwidth, loss behavior, RTT gradient, queue delay, ACK gap, and decoding status. Random loss and congestion-related loss are treated differently: random loss should be absorbed primarily by FEC recovery, whereas loss accompanied by RTT growth and queue buildup should trigger more conservative pacing. By using explicit states rather than an opaque continuous action, the system remains interpretable and easier to debug.

\subsection{Conservative Execution and FEC Reliability Floor}
The execution layer translates the decision state into bounded pacing and FEC actions. Instead of discarding the original PEPesc execution logic, the enhanced system uses a legacy baseline plus adaptive correction strategy. The adaptive correction is limited by step-size and safety constraints, reducing the risk of oscillation caused by prediction errors.

FEC control is also constrained. In dynamic scenarios, reducing redundancy too aggressively may cause decoding stalls, feedback interruption, and connection instability. Therefore, the proposed framework introduces an FEC reliability floor, which ensures that the actual repair-packet ratio does not fall below the minimum level required for streaming recovery. This design treats FEC as a reliability baseline rather than a freely compressible overhead.

\section{Implementation and Experimental Setup}
\label{sec:setup}

\subsection{Prototype Implementation}
The prototype is implemented by extending the PEPesc proxy program. The proxy entities run at the two sides of the controlled bottleneck link and intercept traffic using transparent proxy rules. The enhanced modules are integrated into the proxy-side control loop. The perception module supports the original rule-based mode and the hybrid GRU mode. The control module supports the original legacy mode and the proposed adaptive mode. This separation allows the original PEPesc behavior and the enhanced behavior to be evaluated under the same topology and traffic conditions.

The GRU model is trained offline from logged proxy observations and then used online for inference. The inference output is not directly applied to packet transmission. Instead, it is first combined with the rule-based baseline and then processed by the decision and execution layers. This deployment choice follows the conservative design principle of learning-assisted systems: the learning component improves perception, while the final control remains bounded and interpretable.

\subsection{Emulation Topology}
The evaluation uses a Mininet-based four-node topology. Node A acts as the sender, node D acts as the receiver, and nodes B and C host the two PEP entities. The B--C link is the controlled high-latency bottleneck link. Unless otherwise specified, the bottleneck link is configured with a one-way propagation delay of 300~ms, resulting in a baseline RTT of about 600~ms. Bandwidth and loss rate are controlled by Linux traffic control to create static, dynamic, and long-duration switching scenarios.

Traffic is generated by iperf. The proxy logs runtime variables including estimated bandwidth, RTT, queue delay, ACK timing, loss information, FEC state, and control decisions. These logs are used both for model training and for post-experiment analysis.

\subsection{Evaluation Scenarios and Metrics}
We evaluate three groups of scenarios. Static scenarios keep bandwidth and loss rate unchanged during each run and are used to verify whether the enhancement preserves steady-state capability. Dynamic scenarios switch bandwidth and loss rate across multiple stages and are used to evaluate recovery after link changes. Long-duration dynamic scenarios repeat a base dynamic pattern multiple times and are used to test whether the system can maintain effective operation over a longer period.

The primary metrics include receiver throughput, recovery time after link switching, throughput deficit area, queue delay, RTT, and effective running time. Receiver throughput reflects useful delivery performance. Recovery time measures how quickly the system returns to a stable throughput level after a link change. Throughput deficit area captures the cumulative performance gap during recovery. The p95 queue delay and p95 RTT measure tail delay behavior. Effective running time measures the duration for which the connection maintains usable operation under long-duration switching.

\subsection{Compared Systems}
The baseline system is the original PEPesc configuration using rule-based perception and legacy control. The enhanced system uses hybrid GRU perception and adaptive control with ACCEL/HOLD/BRAKE decisions, conservative pacing correction, and the FEC reliability floor. Both systems are evaluated under the same topology, traffic generation method, and scenario configurations.

\section{Evaluation Results}
\label{sec:evaluation}

\subsection{Static Scenarios}
Static scenarios are used to examine whether the proposed enhancement preserves the steady-state capability of the original proxy system. Across five static scenarios, the enhanced system achieves higher receiver throughput than the original PEPesc baseline, with an average improvement of approximately 7.6\%. This result indicates that the added GRU-assisted perception and adaptive control logic do not introduce a steady-state penalty. Instead, the enhanced perception and bounded pacing correction slightly improve link utilization under stable conditions.

\begin{table}[t]
\centering
\caption{Summary of main experimental results.}
\label{tab:main_results}
\begin{tabular}{lll}
\toprule
Metric & Baseline PEPesc & Enhanced system \\
\midrule
Average static throughput & -- & +7.6\% \\
Average dynamic throughput & 13.17 Mbps & 13.62 Mbps \\
Recovery time after switching & 1.089 s & 0.664 s \\
Throughput deficit area & -- & -20.0\% \\
p95 queue delay & 137.6 ms & 44.1 ms \\
Long-duration effective time & 100.9 s & 188.4 s \\
\bottomrule
\end{tabular}
\end{table}

\subsection{Representative Dynamic Scenario}
The representative dynamic scenario contains multiple stages with changing bandwidth and loss rate. Compared with the baseline, the enhanced system improves average receiver throughput from 13.17~Mbps to 13.62~Mbps. More importantly, it reduces the average recovery time after link switching from 1.089~s to 0.664~s. The throughput deficit area is reduced by approximately 20.0\%, showing that the enhanced system can recover more quickly after link changes.

The improvement is mainly attributed to two factors. First, the hybrid residual GRU reduces the lag of link perception near switching boundaries, allowing the proxy to recognize bandwidth and loss changes earlier. Second, the ACCEL/HOLD/BRAKE mechanism provides explicit recovery and suppression semantics. When the link improves, the system can enter ACCEL and increase pacing conservatively. When queue pressure grows, the system can enter BRAKE and suppress traffic injection before excessive backlog accumulates.

\subsection{Queue Delay and RTT Behavior}
Queue control is a key objective in dynamic high-latency networks. In the representative dynamic scenario, the p95 queue delay decreases from 137.6~ms to 44.1~ms, corresponding to a reduction of about 68.0\%. The p95 RTT also decreases, indicating that the improvement is not limited to average throughput. Lower tail delay suggests that the enhanced system avoids prolonged queue buildup during link transitions.

This behavior is consistent with the execution design. The adaptive correction is bounded rather than aggressive, and the BRAKE state is triggered when queue and RTT signals indicate feedback deterioration. As a result, the system can improve recovery while maintaining stronger queue discipline.

\subsection{Long-Duration Dynamic Switching}
Long-duration experiments evaluate whether the control loop remains stable when dynamic switching continues for an extended period. In these experiments, the baseline system suffers from shorter effective operation under repeated link changes. The enhanced system increases average effective running time from 100.9~s to 188.4~s, an improvement of approximately 86.7\%.

The long-duration result is important because a method that only improves short-term throughput may still fail under continuous dynamics. The observed improvement indicates that the combination of residual perception, state-based decision, conservative execution, and FEC reliability protection improves not only local recovery behavior but also long-term robustness.

\section{Discussion}
\label{sec:discussion}

The evaluation results show that the proposed enhancement improves several aspects of networking service quality in dynamic high-latency environments. The throughput improvement in static scenarios suggests that the additional control logic does not degrade the original steady-state behavior. The reduction in recovery time and throughput deficit area indicates faster adaptation after link changes. The decrease in p95 queue delay shows that the enhanced system also improves tail delay behavior rather than merely increasing throughput. Finally, the long-duration results demonstrate improved robustness under repeated dynamic switching.

The main reason for these improvements is the division of responsibilities among the three layers. The GRU-assisted perception layer improves the timeliness of link-state estimation, but it does not directly control packet transmission. The state-machine decision layer turns noisy estimates and feedback signals into interpretable control semantics. The execution layer then applies bounded corrections and protects the minimum FEC redundancy needed for streaming recovery. This structure avoids the risk of a fully neural control loop while still benefiting from temporal learning.

There are also limitations. First, the evaluation is based on Mininet emulation with controlled link scenarios. Although this setting is useful for repeatable experiments, it cannot cover all dynamics of real satellite or non-terrestrial networks. Second, the GRU model is trained from the available experimental traces, and its generalization to unseen network patterns requires further validation. Third, the current implementation focuses on a single long-lived TCP flow through the proxy. Multi-flow interaction, fairness, encrypted transport compatibility, and deployment issues remain future work.

Despite these limitations, the results indicate that lightweight learning-assisted perception can be useful when it is combined with explicit control semantics and conservative execution constraints. For PEP-based systems, this design direction is more practical than directly replacing the control loop with an opaque model, especially in high-latency environments where unstable actions may persist for one or more RTTs before their consequences become visible.

\section{Conclusion}
\label{sec:conclusion}

This paper presented a GRU-assisted performance enhancing proxy for adaptive transmission over dynamic high-latency networks. Built on PEPesc, the proposed framework integrates hybrid residual GRU-based link perception, ACCEL/HOLD/BRAKE state-based decision making, conservative pacing correction, and an FEC reliability floor. The design preserves the stable rule-based baseline while using temporal learning to improve perception under dynamic bandwidth and loss conditions.

Mininet-based experiments show that the enhanced system improves receiver throughput in static scenarios, accelerates recovery after dynamic link switching, reduces throughput deficit and p95 queue delay, and extends effective running time under long-duration dynamic switching. These results suggest that combining lightweight neural perception with interpretable control states and conservative execution constraints is a practical approach for improving transmission quality in dynamic high-latency PEP systems.

Future work will evaluate the framework under more diverse network conditions, including multi-flow traffic, more complex satellite-like dynamics, and encrypted transport scenarios. Additional work is also needed to study online adaptation and model robustness when the deployment environment differs from the training traces.


\bibliographystyle{splncs04}
\bibliography{references}

\end{document}
